\documentclass[11pt]{article}
\usepackage[a4paper,margin=27mm]{geometry}
\usepackage{amsmath,amssymb,mathtools}
\usepackage{fontspec,unicode-math}
\setmainfont{Latin Modern Roman}
\setmathfont{Latin Modern Math}
\usepackage{xurl}
\usepackage[hidelinks]{hyperref}
\newcommand{\sref}[2]{\hyperlink{source:#1:#2}{[S#2]}}
\newcommand{\cataloguescope}[1]{\paragraph{Mathematical question.}#1}
\hypersetup{pdftitle={Problem statement}}
\begin{document}
\section{\texorpdfstring{For any given integer a > 0, are there infinitely many primes p such that a\textasciicircum{}\{p − 1\} $\equiv$ 1 (mod p\textasciicircum{}\{2\})}{For any given integer a > 0, are there infinitely many primes p such that a\textasciicircum{}\{p − 1\} congruent to 1 (mod p\textasciicircum{}\{2\})}}\label{statement}

\subsection{Definitions and mathematical statement}

\paragraph{Shared notation.}$\mathbb N=\{1,2,\ldots\}$, and $\mathbb Z,\mathbb Q,\mathbb R,\mathbb C$ denote the integers, rationals, reals and complex numbers. Any other conventions are specified locally.


Fix an integer $a>0$. A prime $p$ satisfies the base-$a$ Wieferich condition when $a^{p-1}\equiv1\pmod{p^2}$. The congruence itself excludes primes dividing $a$. For $a=1$ it holds for every prime.
\paragraph{Statement recorded in the supplied source.}
For every fixed integer $a>0$, are there infinitely many primes $p$ with $a^{p-1}\equiv1\pmod{p^2}$? \sref{statement}{2}
\subsection{Short English statement}
Does every positive integer base have infinitely many Wieferich primes? The base one case holds for all primes.
\subsection{Sources}
\begin{description}
\item[\hypertarget{source:statement:1}{S1}] ULAM.AI and the UnsolvedMath Contributors, UnsolvedMath: A Curated Collection of Open Mathematics Problems, catalogue entry. CC BY 4.0. \url{https://huggingface.co/datasets/ulamai/UnsolvedMath}.
\item[\hypertarget{source:statement:2}{S2}] Wikipedia contributors, \emph{List of unsolved problems in mathematics}, revision 1366636767 (2026) Number theory, For any given integer a > 0, are there infinitely many primes p such that a\textasciicircum{}\{p − 1\} $\equiv$ 1 (mod p\textasciicircum{}\{2\}). \url{https://en.wikipedia.org/w/index.php?title=List_of_unsolved_problems_in_mathematics&oldid=1366636767#Number_theory}.
\item[\hypertarget{source:statement:3}{S3}] Jonathan Sondow, \emph{Lerch quotients, Lerch primes, Fermat--Wilson quotients, and the Wieferich-non-Wilson primes $2,3,14771$}, arXiv:1110.3113v5 (2012) Introduction and Section 3.2, pp. 1--2 and 9--10. \url{https://arxiv.org/pdf/1110.3113}.
\item[\hypertarget{source:statement:4}{S4}] OEIS Foundation and contributors, \emph{The On-Line Encyclopedia of Integer Sequences}, A258045: Table T(b, m) of largest exponents k such that for p = prime(m) and base b > 1 the congruence b\textasciicircum{}(p-1) == 1 (mod p\textasciicircum{}k) is satisfied, or 0 if no such k exists, read by antidiagonals (downwards)., official data snapshot \texttt{693b5499c9d2} (2026) Name, definition and comments. \url{https://raw.githubusercontent.com/oeis/oeisdata/693b5499c9d248f570dd2d979e8bd1d7eecf689b/seq/A258/A258045.seq}.
\end{description}
\label{end:statement}
\end{document}
